%%Flux_Aero

selon 3 cas de stabilite : "convection libre" / neutre / autre.
Resolution iterative jusqu'� ce que deux valeurs successives de la
longueur de Monin Obukhov diff�rent de moins de 5 % (seuil choisi).
La convergence du processus it�ratif n'est assur� que pour des conditions
instables.

Cas de la convection libre :
\begin{eqnarray*}
ustar = 0.\\
tstar = -999.\\
vlmon = 0.\\
vhs = 0.\\
\end{eqnarray*}

cas neutre :

\begin{eqnarray*}
ustar = frac{K.dv}{log(zu_2/zu_1)}\\
tstar = 0.\\
vlmon = 9999.0.sign(1.,dt_air)\\
vhs = 0.\\
\end{eqnarray*}

autres cas (stratification neutre, soit vlmon infinie) :

\begin{eqnarray*}
vlmon = 1000000.\\
ustar = frac{K . dv}{log(zu_2/zu_1) - psi_m(zu_2/vlm_1) + psi_m(zu_1/vlm_1)}\\
tstar = frac{K . dt_air}{log(zu_2/zu_1) - psi_h(zu_2/vlm_1) + psi_h(zu_1/vlm_1)}\\
vlmon = frac{tk . ustar^2}{karman . g . tstar}\\
\end{eqnarray*}\\
 
Calcul du flux de chaleur sensible vhs :

Cas instable (avec convergence)

vhs = -roo . cp . ustar . tstar

rho : calcul de la masse volumique d'une particule d'air humide
definie par : pression , temperature , humidite
rho     : reel   : masse en kg/m**3
pres   : reel   : pression em mb
temp   : reel   : temperature en celsius
humi   : reel   : humidite en %

fonction tvir\_4m(pres,temp,humi)
calcul de la temp\'erature virtuelle d'une particule
donnee definie par : pression , temperature , humidite
tvir\_4m : reel  : temp\'erature en degr\'es Celsius

%%%
%% Reste � d�crire l'it�ration. En g�n�ral, c'est
%% sous la forme :
%% soit F(x) dont on chreche les z�ros
%% 1) estimation de X0
%% 2) Calcul de F(Xi)
%% 3) Calcul de F'(Xi)
%% 4) Calcul de Xi+1 par une m�thode
%% (Si Newton, xi+1 = Xi - F(xi)/F'(xi) )
%%  5) si pas convergence, bouclage en 2
%%
Cas stable (sans convergence)

$vhs = {-karman^2 . roo . cp . dv . dt_air}{\log(zu_2/zu_1)^2 / 10}$




avec~:\\
$t_0 = 273.15 $ \\
$karman = 0.4 $ \\
$cp = 1005.0 $ \\
$g = 9.81 $ \\

Dyer 74
